% Chunk: Chapter 6 introduction and first part of \S6.1, Derivation of the Rules.
% Source: Weinberg QFT Vol. I, physical pp. 282--289 (printed pp. 259--266).
% Status: source-reviewed, compile-clean, and render-reviewed; equations (6.1.1)--(6.1.16); figures 6.1--6.2 integrated; one footnote; no uncertainties.

In previous chapters the use of covariant free fields in the construction
of the Hamiltonian density has been motivated by the requirement that
the $S$-matrix satisfy Lorentz invariance and cluster decomposition
conditions. With the Hamiltonian density constructed in this way, it makes
no difference which form of perturbation theory we use to calculate the
$S$-matrix; the results will automatically satisfy these invariance and
clustering conditions in each order in the interaction density. Nevertheless,
there are obvious practical advantages in using a version of perturbation
theory in which the Lorentz invariance and cluster decomposition properties
of the $S$-matrix are kept manifest at every stage in the calculation. This
was not true for the perturbation theory used in the 1930s, now known as
`old-fashioned perturbation theory,' described at the beginning of Section
3.5. The great achievement of Feynman, Schwinger, and Tomonaga in the late
1940s was to develop perturbative techniques for calculating the $S$-matrix,
in which Lorentz invariance and cluster decomposition properties are
transparent throughout. This chapter will outline the diagrammatic
calculational technique first described by Feynman at the Poconos Conference
in 1948. Feynman was led to these diagrammatic rules in part through his
development of a path-integral approach, which will be the subject of Chapter
9. In this chapter, we shall use the approach described by Dyson~\hyperref[ref:6.1]{[1]} in 1949,
which until the 1970s was the basis of almost all analyses of perturbation
theory in quantum field theory, and still provides a particularly transparent
introduction to the Feynman rules.

\subsection{Derivation of the Rules}

Our starting point is a formula for the $S$-matrix, obtained by putting
together the Dyson series \eqref{eq:3.5.10} with expression \eqref{eq:4.2.2} for the
free-particle states:
\begin{align}
&S_{\mathbf p'_1\sigma'_1n'_1;\,\mathbf p'_2\sigma'_2n'_2;\,\ldots,
       \mathbf p_1\sigma_1n_1;\,\mathbf p_2\sigma_2n_2;\,\ldots}
\notag\\
&\quad={}
\sum_{N=0}^{\infty}\frac{(-i)^N}{N!}
\int d^4x_1\cdots d^4x_N\,
\bra{\Omega}\cdots
a_{\mathbf p'_2\sigma'_2n'_2}a_{\mathbf p'_1\sigma'_1n'_1}
\notag\\
&\qquad\cdot
T\!\left\{\mathcal H_I(x_1)\cdots\mathcal H_I(x_N)\right\}
a^\dagger_{\mathbf p_1\sigma_1n_1}
a^\dagger_{\mathbf p_2\sigma_2n_2}\cdots\ket{\Omega}.
\tag{6.1.1}
\label{eq:6.1.1}
\end{align}
As a reminder: $\mathbf p$, $\sigma$, and $n$ label particle momenta, spin,
and species; primes denote labels for particles in the final state;
$\ket{\Omega}$ is the free-particle vacuum state; $a$ and $a^\dagger$ are
annihilation and creation operators; $T$ indicates a time-ordering, which puts
the $\mathcal H_I(x)$ in an order in which the arguments $x^0$ decrease from
left to right; and $\mathcal H_I(x)$ is the interaction Hamiltonian density,
taken as a polynomial in the fields and their adjoints
\begin{equation}
\mathcal H_I(x)=\sum_i g_i\mathcal H_{I,i}(x),
\tag{6.1.2}
\label{eq:6.1.2}
\end{equation}
each term $\mathcal H_{I,i}$ being a product of definite numbers of fields and
field adjoints of each type. The field of a particle of species $n$ that
transforms under a particular representation of the homogeneous Lorentz group
(with or without space inversions) is given by
\begin{align}
\psi_\ell(x)={}&\sum_\sigma(2\pi)^{-3/2}\int d^3p\,
\bigl[u_\ell(\mathbf p,\sigma,n)a_{\mathbf p,\sigma,n}e^{ip\cdot x}
\notag\\
&\hspace{43mm}
+v_\ell(\mathbf p,\sigma,n)a^\dagger_{\mathbf p,\sigma,n^c}e^{-ip\cdot x}
\bigr].
\tag{6.1.3}
\label{eq:6.1.3}
\end{align}
Here $n^c$ denotes the antiparticle of the species $n$, and
$\exp(\pm ip\cdot x)$ is calculated with $p^0$ set equal to
$\sqrt{\mathbf p^2+M_n^2}$. The coefficient functions $u_\ell$ and $v_\ell$
depend on the Lorentz transformation properties of the field and the spin of
the particle it describes; they were calculated in Chapter 5. (For instance,
in the scalar field the $u_\ell$ for a particle of energy $E$ is simply
$(2E)^{-1/2}$, while in a Dirac field $u_\ell$ and $v_\ell$ are the normalized
Dirac spinors introduced in Section 5.5.) The index $\ell$ on the field should
here be understood to indicate the particle type and the representation of the
Lorentz group by which the field transforms, as well as including a running
index labelling the components in this representation. There is no need to
deal separately with interactions that involve derivatives of fields; from
our point of view, the derivative of a field \eqref{eq:6.1.3} is just another field
described by \eqref{eq:6.1.3}, with different $u_\ell$ and $v_\ell$. We will here make a
distinction between some particle species that we arbitrarily call
`particles,' for instance electrons, protons, etc., and those we call
`antiparticles,' such as positrons and antiprotons. The field operators that
destroy particles and create antiparticles are called simply `fields'; their
adjoints, which destroy antiparticles and create particles, are called `field
adjoints.' Of course, some particle species like the photon and $\pi^0$ are
their own antiparticles; for these the field adjoints are proportional to the
fields.

We now proceed to move all annihilation operators to the right in Eq. \eqref{eq:6.1.1},
repeatedly using for this purpose the commutation or anticommutation relations:
\begin{align}
a_{\mathbf p\sigma n}a^\dagger_{\mathbf p'\sigma'n'}
={}&\mathbin{\pm}a^\dagger_{\mathbf p'\sigma'n'}a_{\mathbf p\sigma n}
+\delta^3(\mathbf p'-\mathbf p)\delta_{\sigma'\sigma}\delta_{n'n},
\tag{6.1.4}\label{eq:6.1.4}\\
a_{\mathbf p\sigma n}a_{\mathbf p'\sigma'n'}
={}&\mathbin{\pm}a_{\mathbf p'\sigma'n'}a_{\mathbf p\sigma n},
\tag{6.1.5}\label{eq:6.1.5}\\
a^\dagger_{\mathbf p\sigma n}a^\dagger_{\mathbf p'\sigma'n'}
={}&\mathbin{\pm}a^\dagger_{\mathbf p'\sigma'n'}a^\dagger_{\mathbf p\sigma n}
\tag{6.1.6}
\label{eq:6.1.6}
\end{align}
(and likewise for antiparticles), the $\pm$ sign on the right being $-$ if
both particles $n,n'$ are fermions, and $+$ if either or both are bosons.
Whenever an annihilation operator appears on the extreme right (or a creation
operator on the extreme left), the corresponding contribution to Eq. \eqref{eq:6.1.1}
vanishes, because these operators annihilate the vacuum state:
\begin{align}
a_{\mathbf p\sigma n}\ket{\Omega}&=0,
\tag{6.1.7}\label{eq:6.1.7}\\
\bra{\Omega}a^\dagger_{\mathbf p\sigma n}&=0.
\tag{6.1.8}
\label{eq:6.1.8}
\end{align}
The remaining contributions to Eq. \eqref{eq:6.1.1} are those arising from the delta
function terms on the right-hand side of Eq. \eqref{eq:6.1.4}, with every creation and
annihilation operator in the initial or final states or in the interaction
Hamiltonian density paired in this way with some other annihilation or creation
operator.

In this way, the contribution to Eq. \eqref{eq:6.1.1} of a given order in each of the
terms $\mathcal H_{I,i}$ in the polynomial
$\mathcal H_I(\psi(x),\psi^\dagger(x))$ is given by a sum, over all ways of
pairing creation and annihilation operators~\hyperref[ref:6.2]{[2]}, of the integrals of products
of factors, as follows:

\noindent\textbf{(a)} Pairing of a final particle having quantum numbers
$\mathbf p',\sigma',n'$ with a field adjoint $\psi_\ell^\dagger(x)$ in
$\mathcal H_{I,i}(x)$ yields a factor
\begin{equation}
\left[a_{\mathbf p'\sigma'n'},\psi_\ell^\dagger(x)\right]_{\mp}
=(2\pi)^{-3/2}e^{-ip'\cdot x}u_\ell^*(\mathbf p',\sigma',n').
\tag{6.1.9}
\label{eq:6.1.9}
\end{equation}

\noindent\textbf{(b)} Pairing of a final antiparticle having quantum numbers
$\mathbf p',\sigma',n'^c$ with a field $\psi_\ell(x)$ in
$\mathcal H_{I,i}(x)$ yields a factor
\begin{equation}
\left[a_{\mathbf p'\sigma'n'^c},\psi_\ell(x)\right]_{\mp}
=(2\pi)^{-3/2}e^{-ip'\cdot x}v_\ell(\mathbf p',\sigma',n').
\tag{6.1.10}
\label{eq:6.1.10}
\end{equation}

\noindent\textbf{(c)} Pairing of an initial particle having quantum numbers
$\mathbf p,\sigma,n$ with a field $\psi_\ell(x)$ in
$\mathcal H_{I,i}(x)$ yields a factor
\begin{equation}
\left[\psi_\ell(x),a^\dagger_{\mathbf p\sigma n}\right]_{\mp}
=(2\pi)^{-3/2}e^{ip\cdot x}u_\ell(\mathbf p,\sigma,n).
\tag{6.1.11}
\label{eq:6.1.11}
\end{equation}

\noindent\textbf{(d)} Pairing of an initial antiparticle having quantum
numbers $\mathbf p,\sigma,n^c$ with a field adjoint
$\psi_\ell^\dagger(x)$ in $\mathcal H_{I,i}(x)$ yields a factor
\begin{equation}
\left[\psi_\ell^\dagger(x),a^\dagger_{\mathbf p\sigma n^c}\right]_{\mp}
=(2\pi)^{-3/2}e^{ip\cdot x}v_\ell^*(\mathbf p,\sigma,n).
\tag{6.1.12}
\label{eq:6.1.12}
\end{equation}

\noindent\textbf{(e)} Pairing of a final particle (or antiparticle) having
numbers $\mathbf p',\sigma',n'$ with an initial particle (or antiparticle)
having quantum numbers $\mathbf p,\sigma,n$ yields a factor
\begin{equation}
\left[a_{\mathbf p'\sigma'n'},a^\dagger_{\mathbf p\sigma n}\right]_{\mp}
=\delta^3(\mathbf p'-\mathbf p)\delta_{\sigma'\sigma}\delta_{n'n}.
\tag{6.1.13}
\label{eq:6.1.13}
\end{equation}

\noindent\textbf{(f)} Pairing of a field $\psi_\ell(x)$ in
$\mathcal H_{I,i}(x)$ with a field adjoint $\psi_m^\dagger(y)$ in
$\mathcal H_{I,j}(y)$ yields a factor\footnote{If the interaction
$\mathcal H_I(x)$ is written in the normal-ordered form, as in Eq. \eqref{eq:5.1.33},
then there is no pairing of fields and field adjoints in the \emph{same}
interaction. Otherwise some sort of regularization is needed to give meaning
to $\Delta_{\ell m}(0)$.}
\begin{align}
&\theta(x-y)
\left[\psi_\ell^{(+)}(x),\psi_m^{(+)\dagger}(y)\right]_{\mp}
\mathbin{\pm}\theta(y-x)
\left[\psi_m^{(-)\dagger}(y),\psi_\ell^{(-)}(x)\right]_{\mp}
\notag\\
&\hspace{48mm}\equiv-i\Delta_{\ell m}(x,y),
\tag{6.1.14}
\label{eq:6.1.14}
\end{align}
where $\psi^{(+)}$ and $\psi^{(-)}$ are the terms in $\psi$ that destroy
particles and create antiparticles, respectively:
\begin{align}
\psi_\ell^{(+)}(x)
&=(2\pi)^{-3/2}\int d^3p\sum_\sigma
u_\ell(\mathbf p\sigma n)e^{ip\cdot x}a_{\mathbf p\sigma n},
\tag{6.1.15}\label{eq:6.1.15}\\
\psi_\ell^{(-)}(x)
&=(2\pi)^{-3/2}\int d^3p\sum_\sigma
v_\ell(\mathbf p\sigma n)e^{-ip\cdot x}a^\dagger_{\mathbf p\sigma n^c}.
\tag{6.1.16}
\label{eq:6.1.16}
\end{align}
Recall that $\theta(x-y)$ is a step function, equal to $+1$ for $x^0>y^0$
and zero for $x^0<y^0$. These step functions appear in Eq. \eqref{eq:6.1.14} because
of the time-ordering in Eq. \eqref{eq:6.1.1}; we can encounter a pairing of an
annihilation field $\psi^{(+)}(x)$ in $\mathcal H_I(x)$ with a creation field
$\psi^{(+)\dagger}(y)$ in $\mathcal H_I(y)$ only if $\mathcal H_I(x)$ was initially
to the left of $\mathcal H_I(y)$ in Eq. \eqref{eq:6.1.1}, i.e., if $x^0>y^0$;
similarly, we encounter a pairing of an annihilation field
$\psi^{(-)\dagger}(y)$ in $\mathcal H_I(y)$ with a creation field
$\psi^{(-)}(x)$ in $\mathcal H_I(x)$ only if $\mathcal H_I(y)$ was initially
to the left of $\mathcal H_I(x)$ in Eq. \eqref{eq:6.1.1}, i.e., if $y^0>x^0$. (The
$\pm$ sign in the second term in \eqref{eq:6.1.14} will be explained a little later.)
The quantity \eqref{eq:6.1.14} is known as a \emph{propagator}; it is calculated in the
following section.

The $S$-matrix is obtained by multiplying these factors together, along with
additional numerical factors to be discussed below, then integrating over
$x_1\cdots x_N$, then summing over all pairings, and then over the numbers of
interaction of each type. Before filling in all the details, it will be
convenient first to describe a diagrammatic formalism for keeping track of all
these pairings.

\begingroup
\renewcommand{\thefigure}{6.\arabic{figure}}
\setcounter{figure}{0}
\begin{figure}[htbp]
% Figure 6.1 diagram plate only; caption remains in the chapter text.
\begin{center}
\begin{tikzpicture}[
  x=1cm,
  y=1cm,
  line cap=round,
  line join=round,
  every node/.style={font=\small},
  particle/.style={line width=0.75pt},
  flow/.style={-{Stealth[length=2.4mm,width=1.8mm]},line width=0.75pt}
]
% (a) Final particle paired with a field adjoint.
\draw[particle] (-6.25,6.55) -- (-6.25,4.05);
\draw[flow] (-6.25,5.02) -- (-6.25,5.48);
\fill (-6.25,4.05) circle[radius=1.7pt];
\draw[particle] (-6.25,4.05) -- (-6.47,3.75);
\draw[particle] (-6.25,4.05) -- (-6.03,3.75);
\node[anchor=south] at (-6.25,6.62)
  {$\mathbf p'\sigma'n'$};
\node[anchor=north west] at (-6.08,3.98) {$\ell,x$};
\node at (-4.35,5.25)
  {$\displaystyle
    \frac{e^{-\ii p'\cdot x}u_\ell^*(\mathbf p'\sigma'n')}
         {(2\pi)^{3/2}}$};
\node at (-4.70,3.45) {(a)};

% (b) Final antiparticle paired with a field.
\draw[particle] (-1.35,6.55) -- (-1.35,4.05);
\draw[flow] (-1.35,5.48) -- (-1.35,5.02);
\fill (-1.35,4.05) circle[radius=1.7pt];
\draw[particle] (-1.35,4.05) -- (-1.57,3.75);
\draw[particle] (-1.35,4.05) -- (-1.13,3.75);
\node[anchor=south] at (-1.35,6.62)
  {$\mathbf p'\sigma'n'^c$};
\node[anchor=north west] at (-1.18,3.98) {$\ell,x$};
\node at (0.55,5.25)
  {$\displaystyle
    \frac{e^{-\ii p'\cdot x}v_\ell(\mathbf p'\sigma'n')}
         {(2\pi)^{3/2}}$};
\node at (0.20,3.45) {(b)};

% (c) Initial particle paired with a field.
\draw[particle] (3.55,3.78) -- (3.55,6.28);
\draw[flow] (3.55,4.82) -- (3.55,5.28);
\fill (3.55,6.28) circle[radius=1.7pt];
\draw[particle] (3.55,6.28) -- (3.33,6.58);
\draw[particle] (3.55,6.28) -- (3.77,6.58);
\node[anchor=south west] at (3.69,6.31) {$\ell,x$};
\node[anchor=north] at (3.55,3.70)
  {$\mathbf p\sigma n$};
\node at (5.45,5.25)
  {$\displaystyle
    \frac{e^{\ii p\cdot x}u_\ell(\mathbf p\sigma n)}
         {(2\pi)^{3/2}}$};
\node at (5.10,3.45) {(c)};

% (d) Initial antiparticle paired with a field adjoint.
\draw[particle] (-6.25,-0.02) -- (-6.25,2.48);
\draw[flow] (-6.25,1.45) -- (-6.25,0.99);
\fill (-6.25,2.48) circle[radius=1.7pt];
\draw[particle] (-6.25,2.48) -- (-6.47,2.78);
\draw[particle] (-6.25,2.48) -- (-6.03,2.78);
\node[anchor=south west] at (-6.11,2.51) {$\ell,x$};
\node[anchor=north] at (-6.25,-0.10)
  {$\mathbf p\sigma n^c$};
\node at (-4.35,1.23)
  {$\displaystyle
    \frac{e^{\ii p\cdot x}v_\ell^*(\mathbf p\sigma n)}
         {(2\pi)^{3/2}}$};
\node at (-4.70,-0.45) {(d)};

% (e) Final particle or antiparticle paired directly with the initial one.
\draw[particle] (-0.95,-0.02) -- (-0.95,2.60);
\draw[flow] (-0.95,1.05) -- (-0.95,1.51);
\node[anchor=south] at (-0.95,2.67)
  {$\mathbf p'\sigma'n'$};
\node[anchor=north] at (-0.95,-0.10)
  {$\mathbf p\sigma n$};

\node at (0.18,1.28) {or};

\draw[particle] (1.28,-0.02) -- (1.28,2.60);
\draw[flow] (1.28,1.51) -- (1.28,1.05);
\node[anchor=south] at (1.28,2.67)
  {$\mathbf p'\sigma'n'^c$};
\node[anchor=north] at (1.28,-0.10)
  {$\mathbf p\sigma n^c$};
\node[anchor=west] at (2.18,1.28)
  {$\delta^3(\mathbf p'-\mathbf p)
    \delta_{\sigma'\sigma}\delta_{n'n}$};
\node at (0.18,-0.55) {(e)};

% (f) Field paired with a field adjoint.
\draw[particle] (-2.15,-2.05) -- (1.00,-2.05);
\draw[flow] (-0.28,-2.05) -- (-0.74,-2.05);
\fill (-2.15,-2.05) circle[radius=1.7pt];
\fill (1.00,-2.05) circle[radius=1.7pt];
\draw[particle] (-2.15,-2.05) -- (-2.48,-1.83);
\draw[particle] (-2.15,-2.05) -- (-2.48,-2.27);
\draw[particle] (1.00,-2.05) -- (1.33,-1.83);
\draw[particle] (1.00,-2.05) -- (1.33,-2.27);
\node[anchor=north] at (-2.15,-2.22) {$\ell,x$};
\node[anchor=north] at (1.00,-2.22) {$m,y$};
\node[anchor=west] at (1.95,-2.05)
  {$-\ii\Delta_{\ell m}(x,y)$};
\node at (-0.58,-3.02) {(f)};
\end{tikzpicture}
\end{center}

\caption{Graphical representation of pairings of operators arising in the
coordinate-space evaluation of the $S$-matrix. The expressions on the right
are the factors that must be included in the coordinate-space integrand of the
$S$-matrix for each line of the Feynman diagram.}
\label{fig:6.1}
\end{figure}
\endgroup

The rules for calculating the $S$-matrix are conveniently summarized in terms
of \emph{Feynman diagrams}. (See Figure~\ref{fig:6.1}.) The diagrams consist of points
called \emph{vertices}, each representing one of the $\mathcal H_{I,i}(x)$, and
\emph{lines}, each representing the pairing of a creation with an annihilation
operator. More specifically:

\noindent\textbf{(a)} The pairing of a final particle with a field adjoint in
one of the $\mathcal H_{I,i}(x)$ is represented by a line running from the vertex
representing that $\mathcal H_{I,i}(x)$ upwards out of the diagram, carrying an
arrow pointed upwards.

\noindent\textbf{(b)} The pairing of a final antiparticle with a field in one
of the $\mathcal H_{I,i}(x)$ is also represented by a line running from the vertex
representing that $\mathcal H_{I,i}(x)$ upwards out of the diagram, but carrying
an arrow pointed downwards. (Arrows are omitted throughout for particles like
$\gamma$, $\pi^0$, etc. that are their own antiparticles.)

\noindent\textbf{(c)} The pairing of an initial particle with a field in one
of the $\mathcal H_{I,i}(x)$ is represented by a line running into the diagram from
below, ending in the vertex representing that $\mathcal H_{I,i}(x)$, carrying an
arrow pointed upwards.

\noindent\textbf{(d)} The pairing of an initial antiparticle with a field in
one of the $\mathcal H_{I,i}(x)$ is also represented by a line running into the
diagram from below, ending in the vertex representing that $\mathcal H_{I,i}(x)$,
but carrying an arrow pointed downwards.

\noindent\textbf{(e)} The pairing of a final particle or antiparticle with an
initial particle or antiparticle is represented by a line running clear
through the diagram from bottom to top, not touching any vertex, with an arrow
pointed upwards or downwards for particles or antiparticles, respectively.

\noindent\textbf{(f)} The pairing of a field in $\mathcal H_{I,i}(x)$ with a
field adjoint in $\mathcal H_{I,j}(y)$ is represented by a line joining the
vertices representing $\mathcal H_{I,i}(x)$ and $\mathcal H_{I,j}(y)$,
carrying an arrow pointing from $y$ to $x$.

Note that arrows always point in the direction a particle is moving, and
opposite to the direction an antiparticle is moving. (As mentioned above,
arrows should be omitted for particles like photons that are their own
antiparticles.) The arrow direction indicated in rule (f) is consistent with
this convention because a field adjoint in $\mathcal H_{I,j}(y)$ can either
create a particle destroyed by a field in $\mathcal H_{I,i}(x)$, or destroy an
antiparticle created by a field in $\mathcal H_{I,i}(x)$. Note also that since
every field or field adjoint in $\mathcal H_{I,i}(x)$ must be paired with
something, the total number of lines at a vertex of type $i$, corresponding to
a term $\mathcal H_{I,i}(x)$ in Eq. \eqref{eq:6.1.2}, is just equal to the total number
of field or field adjoint factors in $\mathcal H_{I,i}(x)$. Of these lines, the
number with arrows pointed into the vertex or out of it equals the number of
fields or field adjoints respectively in the corresponding interaction term.

To calculate the contribution to the $S$-matrix for a given process, of a given
order $N_i$ in each of the interaction terms $\mathcal H_{I,i}(x)$ in Eq.
\eqref{eq:6.1.2}, we must carry out the following steps:

\noindent\textbf{(i)} Draw all Feynman diagrams containing $N_i$ vertices of
each type $i$, and containing a line coming into the diagrams from below for
each particle or antiparticle in the initial state, and a line going upwards
out of the diagram for every particle or antiparticle in the final state,
together with any number of internal lines running from one vertex to another,
as required to give each vertex the proper number of attached lines. The lines
carry arrows as described above, each of which may point upwards or downwards.
Each vertex is labelled with an interaction type $i$ and spacetime coordinate
$x^\mu$. Each internal or external line is labelled at the end where it runs
into a vertex with a field type $\ell$ (corresponding to the field
$\psi_\ell(x)$ or $\psi_\ell^\dagger(x)$ that creates or destroys the particle
or antiparticle at that vertex), and each external line where it enters or
leaves the diagram is labelled with the quantum numbers $\mathbf p,\sigma,n$
or $\mathbf p',\sigma',n'$ of the initial or final particle (or antiparticle).

\noindent\textbf{(ii)} For each vertex of type $i$, include a factor $-i$
(from the $(-i)^N$ in Eq. \eqref{eq:6.1.1}) and a factor $g_i$ (the coupling constant
multiplying the product of fields in $\mathcal H_{I,i}(x)$). For each line
running upwards out of the diagram, include a factor \eqref{eq:6.1.9} or \eqref{eq:6.1.10},
depending on whether the arrow is pointing up or down. For each line running
from below into the diagram, include a factor \eqref{eq:6.1.11} or \eqref{eq:6.1.12}, again
depending on the arrow direction. For each line running straight through the
diagram include a factor \eqref{eq:6.1.13}. For each internal line connecting two
vertices include a factor \eqref{eq:6.1.14}.

\noindent\textbf{(iii)} Integrate the product of all these factors over the
coordinates $x_1,x_2,\ldots$ of each vertex.

\noindent\textbf{(iv)} Add up the results obtained in this way from each
Feynman diagram. The complete perturbation series for the $S$-matrix is
obtained by adding up the contributions of each order in each interaction
type, up to whatever order our strength permits.

Note that we have not included the factor $1/N!$ from Eq. \eqref{eq:6.1.1} in these
rules, because the time-ordered product in Eq. \eqref{eq:6.1.1} is a sum over the $N!$
permutations of $x_1x_2\cdots x_N$, each permutation giving the same
contribution to the final result. To put this another way, a Feynman diagram
with $N$ vertices is one of $N!$ identical diagrams, which differ only in
permutations of the labels on the vertices, and this yields a factor of $N!$
which cancels the $1/N!$ in Eq. \eqref{eq:6.1.1}. (There are exceptions to this rule,
discussed below.) For this reason, henceforth we do \emph{not} include more
than one of a set of Feynman diagrams that differ only by relabelling the
vertices.

In some cases, there are additional combinatoric factors or signs that must be
included in the contribution of individual Feynman diagrams:

\noindent\textbf{(v)} Suppose that an interaction $\mathcal H_{I,i}(x)$
contains (among other fields and field adjoints) $M$ factors of the same field.
Suppose that each of these fields is paired with a field adjoint in a different
interaction (different for each one), or in the initial or final state. The
first of these field adjoints
\begingroup
\renewcommand{\thefigure}{6.\arabic{figure}}
\setcounter{figure}{1}
\begin{figure}[htbp]
% Figure 6.2 diagram only; caption remains in the chapter text.
\begin{center}
\begin{tikzpicture}[
  x=1cm,
  y=1cm,
  line cap=round,
  line join=round,
  particle/.style={line width=0.8pt},
  flow/.style={-{Stealth[length=2.5mm,width=1.9mm]},line width=0.8pt}
]
\coordinate (L) at (-1.65,0);
\coordinate (R) at (1.65,0);

\draw[particle] (-3.15,0) -- (L);
\draw[particle] (R) -- (3.15,0);
\draw[particle] (L) -- (R);
\draw[flow] (-0.24,0) -- (0.28,0);

\draw[particle] (L) .. controls (-1.20,1.42) and (1.20,1.42) .. (R);
\draw[flow] (-0.26,1.065) -- (0.27,1.065);

\draw[particle] (L) .. controls (-1.20,-1.42) and (1.20,-1.42) .. (R);
\draw[flow] (-0.26,-1.065) -- (0.27,-1.065);

\fill (L) circle[radius=2pt];
\fill (R) circle[radius=2pt];
\end{tikzpicture}
\end{center}

\caption{Example of a graph requiring extra combinatoric factors in the
$S$-matrix. For an interaction involving, say, three factors of some field (as
well as other fields) we usually include a factor $1/3!$ in the interaction
Hamiltonian, to cancel factors arising from sums over ways of pairing these
fields with their adjoints in other interactions. But in this diagram there
are two such factors of $1/3!$ and only $3!$ different pairings, so we are left
with an extra factor of $1/3!$.}
\label{fig:6.2}
\end{figure}
\endgroup

\noindent can be paired with any one of the $M$ identical fields in
$\mathcal H_{I,i}(x)$; the second with any one of the remaining $M-1$
identical fields; and so on, yielding an extra factor of $M!$. To compensate
for this, it is conventional to define the coupling constants $g_i$ so that an
explicit factor $1/M!$ appears in any $\mathcal H_{I,i}(x)$ containing $M$
identical fields (or field adjoints.) For instance, the interaction of $M$th
order in a scalar field $\phi(x)$ would be written $g\phi^M/M!$. (More
generally, one often also displays an explicit factor of $1/M!$ when the
interaction involves a sum of $M$ factors of fields from the same symmetry
multiplet, or when for this or any other reason the coupling coefficient is
totally symmetric or antisymmetric under permutations of $M$ boson or fermion
fields.)

However, this cancellation of $M!$ factors is not always complete. For
instance, consider a Feynman diagram in which the $M$ identical fields in one
interaction $\mathcal H_{I,i}(x)$ are paired with $M$ corresponding field
adjoints in a \emph{single} other interaction $\mathcal H_{I,j}(y)$. (See
Figure~\ref{fig:6.2}.) Then by following the above analysis, we find only $M!$ different
pairings (since it makes no difference which of the field adjoints we call the
first, second, $\ldots$), cancelling only one of the two factors of $1/M!$ in
the two different interactions. In this case, we would have to insert an extra
factor of $1/M!$ `by hand' into the contribution of such a Feynman diagram.
